\section{Introduction}
\IEEEPARstart{A}{world} model predicts how an environment evolves in response to an agent's actions. In robot-assisted surgery such a model could act as a photorealistic simulator for training and planning, as an automatic evaluator of autonomous policies \cite{cosmossurgdvrk}, and as a generator of paired video--action data for learning those policies \cite{cosmoshsurg,surgwam}. Medical world models have begun to show this value outside surgery, for example by simulating tumour response to candidate treatments and choosing among them \cite{mewm}.

Building surgical world models is limited by the action signal. Models conditioned on robot kinematics \cite{openh,kvlr,surgcast,surgwam} can be precisely controlled, but synchronized kinematics are recorded only on research platforms such as the da Vinci Research Kit, mostly on phantoms or ex-vivo tissue, totalling a few hundred hours across all public datasets \cite{openh}. Clinical and in-vivo robotic video is far more abundant but carries no kinematics. Action-free approaches learn latent actions from such video in the manner of Genie \cite{genie}: the surgical vision world model of Koju et al.\ \cite{surgwm} infers one of twelve discrete latent actions per frame transition. A single, monolithic latent action, however, has to explain everything that changes between two endoscopic frames---endoscope motion, the motion of two or three instruments, and the tissue's response---and latent-action learning is known to degrade sharply under such action-correlated distractors unless a small amount of supervision is added \cite{laom}.

We start from an observation about teleoperated surgical robots: \emph{their embodiment is observable even when their kinematics are not}. The da Vinci system logs which of its four arms holds the endoscope and which instrument is mounted on every other arm, with timestamps; the console overlay burned into the recorded video shows which arms the surgeon's hands currently drive and when energy is applied; and teleoperation itself is structured, because at most two instrument arms move at once and the endoscope moves in a separate control mode. These signals come for free with thousands of hours of video, yet no surgical world model uses them.

We propose \method{}, an embodiment-grounded surgical world model whose action space mirrors the robot (Fig.~\ref{fig:arch}). The endoscope action is explicit: a four-parameter similarity transform measured from the video itself. Each instrument arm receives its own latent action, inferred by an inverse-dynamics model that observes only that arm's instrument region---located by an instrument detector and bound to the arm through the robot log---and constrained by an information bottleneck. A latent spatio-temporal diffusion transformer predicts future video tokens from context frames, the endoscope action through adaptive normalization, and one token per arm through cross-attention. Because each action is tied to a physical actuator, the model can be steered one arm at a time and its actions can be compared with real kinematics.

Our contributions are:
\begin{itemize}
\item \method{}, to our knowledge the first surgical world model with an embodiment-factorized action space learned without kinematic labels, combining an explicit endoscope ego-motion action with per-arm latent actions grounded in robot system logs.
\item A data pipeline that turns \todo{838}\,h of in-vivo da Vinci video and its system logs into model-ready windows with robot state, endoscope ego-motion and arm-bound instrument regions, and an analysis of the label noise in such logs.
\item An evaluation protocol for surgical world models that goes beyond pixel fidelity: controllability, per-arm locality, compliance with commanded endoscope motion, consistency of the predicted scene with the robot's state, and linear probing of learned actions against real kinematics.
\item Extensive experiments against monolithic \cite{genie,surgwm} and factored \cite{flam} latent actions at equal capacity, kinematics-conditioned baselines \cite{openh} on public benchmarks, and ablations of every component, showing \todo{(summary of findings)}.
\end{itemize}
